% Source: https://artofproblemsolving.com/community/c3h1985346_2019_amc_10b_18
% Stable identity: aops-topic-1985346-q01
Measure positions in kilometers from home, so home is $0$ and the gym is $2$. Let $p_0=0$ and let $p_1,p_2,\dots$ be the points where Henry turns. A walk toward the gym from $p$ ends at $p+\dfrac{3}{4}(2-p)=\dfrac{p}{4}+\dfrac{3}{2}$, and a walk toward home from $p$ ends at $p-\dfrac{3}{4}p=\dfrac {p}{4}$.

So the odd turning points $p_1,p_3,\dots$ are near the gym and the even ones $p_2,p_4,\dots$ are near home, with
\[
p_{2k+1}=\dfrac{p_{2k}}{4}+\dfrac{3}{2},\qquad p_{2k+2}=\dfrac{p_{2k+1}}{4}.
\]
Then $p_{2k+3}=\dfrac{p_{2k+1}}{16}+\dfrac{3}{2}$. The fixed point of $x\mapsto\dfrac{x}{16}+\dfrac{3}{2}$ is $B=\dfrac{8}{5}$, and $p_{2k+3}-B=\dfrac{1}{16}(p_{2k+1}-B)$, so the turning points near the gym approach $B=\dfrac{8}{5}$. The turning points near home are a quarter of these, so they approach $A=\dfrac {B}{4}=\dfrac{2}{5}$. (Check: $\dfrac{2}{5}+\dfrac{3}{4}\left(2-\dfrac{2}{5}\right)=\dfrac{8}{5}$ and $\dfrac{8}{5}-\dfrac{3}{4}\cdot\dfrac{8}{5}=\dfrac{2}{5}$.)

These $A$ and $B$ are the two points between which Henry eventually walks back and forth, and
\[
|A-B|=\dfrac{8}{5}-\dfrac{2}{5}=\boxed{\dfrac{6}{5}}.
\]
